% ============================================================
%  CEFUSA E-Commerce — Explicación detallada del código
%  Arquitecturas de Software 2026-1
% ============================================================
\documentclass[12pt, a4paper]{article}

% ─── Paquetes ────────────────────────────────────────────────
\usepackage[utf8]{inputenc}
\usepackage[T1]{fontenc}
\usepackage[spanish]{babel}
\usepackage{geometry}
\geometry{margin=2.5cm}

\usepackage{xcolor}
\usepackage{listings}
\usepackage{titlesec}
\usepackage{parskip}
\usepackage{hyperref}
\usepackage{booktabs}
\usepackage{tikz}
\usetikzlibrary{shapes.geometric, arrows.meta, positioning, fit, backgrounds}
\usepackage{mdframed}
\usepackage{enumitem}
\usepackage{fancyhdr}
\usepackage{graphicx}
\usepackage{tcolorbox}
\tcbuselibrary{breakable, skins}

% ─── Colores ─────────────────────────────────────────────────
\definecolor{bg}{HTML}{1E1E2E}
\definecolor{fg}{HTML}{CDD6F4}
\definecolor{kw}{HTML}{CBA6F7}
\definecolor{str}{HTML}{A6E3A1}
\definecolor{cmt}{HTML}{6C7086}
\definecolor{num}{HTML}{FAB387}
\definecolor{cls}{HTML}{89DCEB}
\definecolor{accent}{HTML}{89B4FA}
\definecolor{warn}{HTML}{F38BA8}
\definecolor{tip}{HTML}{A6E3A1}
\definecolor{boxbg}{HTML}{F0F4FF}
\definecolor{boxborder}{HTML}{89B4FA}

% ─── Estilo de código Python ─────────────────────────────────
\lstdefinestyle{python}{
  backgroundcolor=\color{bg},
  basicstyle=\ttfamily\small\color{fg},
  keywordstyle=\color{kw}\bfseries,
  stringstyle=\color{str},
  commentstyle=\color{cmt}\itshape,
  numberstyle=\tiny\color{cmt},
  numbers=left,
  numbersep=8pt,
  showstringspaces=false,
  breaklines=true,
  frame=single,
  rulecolor=\color{cmt},
  tabsize=4,
  morekeywords={self, None, True, False, from, import, class, def,
                return, raise, try, except, if, else, elif, for, in,
                not, and, or, pass, with},
  literate=
    {'}{{\textquotesingle}}1
    {á}{{\'a}}1 {é}{{\'e}}1 {í}{{\'i}}1 {ó}{{\'o}}1 {ú}{{\'u}}1
    {Á}{{\'A}}1 {É}{{\'E}}1 {Í}{{\'I}}1 {Ó}{{\'O}}1 {Ú}{{\'U}}1
    {ñ}{{\~n}}1 {Ñ}{{\~N}}1
}
\lstset{style=python}

% ─── Cajas informativas ───────────────────────────────────────
\newtcolorbox{conceptbox}[1]{
  breakable, enhanced,
  colback=boxbg, colframe=boxborder,
  title=\textbf{#1},
  fonttitle=\color{white},
  coltitle=white,
  attach boxed title to top left={yshift=-2mm, xshift=4mm},
  boxed title style={colback=accent},
  sharp corners, boxrule=0.8pt
}

\newtcolorbox{warnbox}[1]{
  breakable, enhanced,
  colback=white!95!red, colframe=warn,
  title=\textbf{#1},
  fonttitle=\color{white},
  coltitle=white,
  attach boxed title to top left={yshift=-2mm, xshift=4mm},
  boxed title style={colback=warn},
  sharp corners, boxrule=0.8pt
}

\newtcolorbox{tipbox}[1]{
  breakable, enhanced,
  colback=white!95!green, colframe=tip,
  title=\textbf{#1},
  fonttitle=\color{white},
  coltitle=white,
  attach boxed title to top left={yshift=-2mm, xshift=4mm},
  boxed title style={colback=tip!70!black},
  sharp corners, boxrule=0.8pt
}

% ─── Cabecera ─────────────────────────────────────────────────
\pagestyle{fancy}
\fancyhf{}
\rhead{\small CEFUSA E-Commerce}
\lhead{\small Arquitecturas de Software 2026-1}
\cfoot{\thepage}

% ─── Titulación ───────────────────────────────────────────────
\titleformat{\section}{\Large\bfseries\color{accent}}{}{0em}{}[\titlerule]
\titleformat{\subsection}{\large\bfseries\color{kw}}{}{0em}{}
\titleformat{\subsubsection}{\normalsize\bfseries\color{cls}}{}{0em}{}

% ═══════════════════════════════════════════════════════════════
\begin{document}
% ═══════════════════════════════════════════════════════════════

\begin{titlepage}
  \centering
  \vspace*{2cm}
  {\Huge\bfseries\color{accent} CEFUSA E-Commerce\par}
  \vspace{0.5cm}
  {\Large Explicación Técnica del Código Fuente\par}
  \vspace{1cm}
  \rule{\linewidth}{1pt}
  \vspace{0.5cm}
  {\large Arquitecturas de Software 2026-1\par}
  {\large Entregable No.\,1 — Núcleo de Negocio y API REST\par}
  \vspace{1.5cm}
  {\normalsize 5 de marzo de 2026\par}
  \vfill
  \textit{Documento autogenerado a partir del código fuente del proyecto}
\end{titlepage}

\tableofcontents
\clearpage

% ═══════════════════════════════════════════════════════════════
\section{Visión General: La Arquitectura en Capas}
% ═══════════════════════════════════════════════════════════════

El proyecto está dividido en \textbf{cuatro capas} bien diferenciadas. Cada capa tiene
una responsabilidad específica y nunca mezcla responsabilidades de otra.

\begin{center}
\begin{tikzpicture}[
  layer/.style={draw=accent, fill=boxbg, rounded corners=4pt,
                minimum width=9cm, minimum height=1cm,
                text width=8.5cm, align=center, font=\small},
  arrow/.style={-Stealth, thick, color=accent}
]
  \node[layer, fill=warn!15] (pres)
    {\textbf{Capa de Presentación} \\ \texttt{views.py} + \texttt{serializers.py} \\ Solo HTTP: recibe, valida y responde};
  \node[layer, fill=kw!10, below=0.4cm of pres] (app)
    {\textbf{Capa de Aplicación} \\ \texttt{services.py} \\ Lógica de negocio y orquestación};
  \node[layer, fill=accent!15, below=0.4cm of app] (dom)
    {\textbf{Capa de Dominio} \\ \texttt{models.py} + \texttt{domain/builders.py} \\ Entidades y construcción};
  \node[layer, fill=str!15, below=0.4cm of dom] (infra)
    {\textbf{Infraestructura} \\ \texttt{infra/factories.py}, \texttt{notifiers.py}, \texttt{payment.py} \\ Dependencias externas};

  \draw[arrow] (pres) -- node[right]{\tiny delega} (app);
  \draw[arrow] (app)  -- node[right]{\tiny usa}     (dom);
  \draw[arrow] (app)  -- node[left]{\tiny instancia via Factory} (infra);
\end{tikzpicture}
\end{center}

\begin{conceptbox}{Regla de oro del proyecto}
\textbf{Una vista nunca toca la base de datos directamente.} Toda lectura o escritura
pasa por un \textit{Service}. Los \textit{Services} usan \textit{Builders} para construir
objetos complejos y \textit{Factories} para obtener dependencias externas.
\end{conceptbox}

% ═══════════════════════════════════════════════════════════════
\section{Capa de Dominio: los Modelos}
% ═══════════════════════════════════════════════════════════════

Los modelos son la base de datos del sistema. En Django, cada clase Python se convierte
en una tabla. Aquí están las \textbf{6 entidades} del proyecto (el 60\% del dominio).

\subsection{Product y ProductVariant}

\begin{lstlisting}[caption={products/models.py -- Producto y Variante}]
class Product(models.Model):
    CATEGORY_CHOICES = [
        ('clothes',    'Ropa'),
        ('accesories', 'Accesorios'),
        ('other',      'Otros')
    ]
    name        = models.CharField(max_length=150)
    description = models.TextField()
    category    = models.CharField(max_length=50, choices=CATEGORY_CHOICES,
                                   default='other')
    is_active   = models.BooleanField(default=True)
    created_at  = models.DateTimeField(auto_now_add=True)
\end{lstlisting}

\begin{conceptbox}{¿Qué es \texttt{choices}?}
El parámetro \texttt{choices} limita los valores aceptados en la base de datos a una
lista predefinida. Django la valida automáticamente y genera etiquetas legibles
(\textit{"Ropa"}, \textit{"Accesorios"}) para la interfaz de administración.
\end{conceptbox}

\begin{lstlisting}[caption={products/models.py -- ProductVariant}]
class ProductVariant(models.Model):
    product      = models.ForeignKey(Product, related_name='variants',
                                     on_delete=models.CASCADE)
    sku          = models.CharField(max_length=100, unique=True)
    size         = models.CharField(max_length=20, null=True, blank=True)
    color        = models.CharField(max_length=30, null=True, blank=True)
    price        = models.DecimalField(max_digits=10, decimal_places=2,
                                       validators=[MinValueValidator(0)])
    is_available = models.BooleanField(default=True)
\end{lstlisting}

\begin{conceptbox}{¿Qué es \texttt{ForeignKey}?}
Una \texttt{ForeignKey} crea una relación \textbf{muchos-a-uno}. Aquí, muchas variantes
pueden pertenecer a un solo producto. \texttt{related\_name='variants'} permite
acceder desde un producto a sus variantes con \texttt{producto.variants.all()}.
\texttt{on\_delete=CASCADE} significa que si se borra el producto, se borran todas
sus variantes automáticamente.
\end{conceptbox}

\subsection{Inventory}

\begin{lstlisting}[caption={products/models.py -- Inventory con métodos de validación}]
class Inventory(models.Model):
    product_variant    = models.OneToOneField(ProductVariant,
                             related_name='inventory', on_delete=models.CASCADE)
    available_quantity = models.PositiveIntegerField(default=0)
    minimum_stock      = models.PositiveIntegerField(default=5)
    updated_at         = models.DateTimeField(auto_now=True)

    def has_stock(self, quantity: int) -> bool:
        """Validacion de negocio: hay suficiente stock?"""
        return self.available_quantity >= quantity

    def is_below_minimum(self) -> bool:
        """Alerta de reposicion: el stock esta bajo el minimo?"""
        return self.available_quantity <= self.minimum_stock
\end{lstlisting}

\begin{conceptbox}{¿Por qué Inventory es una entidad separada?}
El \texttt{Inventory} está separado de \texttt{ProductVariant} intencionalmente.
Esto sigue el principio de \textbf{Separación de Responsabilidades (SRP)}: la variante
describe \textit{qué es} el producto, mientras que el inventario describe
\textit{cuántos hay}. La separación facilita auditorías y permite que los dos
evolucionen independientemente.
\end{conceptbox}

\begin{tipbox}{OneToOneField vs ForeignKey}
\texttt{OneToOneField} es como \texttt{ForeignKey} pero garantiza que \textbf{un solo
inventario} corresponde a \textbf{una sola variante}. Si la variante es la llave
del cuarto, el inventario es el único baúl que puede tener esa llave.
\end{tipbox}

\subsection{Customer}

\begin{lstlisting}[caption={customers/models.py -- Cliente}]
class Customer(models.Model):
    nombre     = models.CharField(max_length=100)
    apellido   = models.CharField(max_length=100)
    email      = models.EmailField(unique=True)
    telefono   = models.CharField(max_length=20, blank=True)
    direccion  = models.TextField(blank=True)
    created_at = models.DateTimeField(auto_now_add=True)

    @property
    def nombre_completo(self):
        return f"{self.nombre} {self.apellido}"
\end{lstlisting}

\begin{conceptbox}{¿Qué es \texttt{@property}?}
Un \texttt{@property} permite acceder a \texttt{customer.nombre\_completo} como si
fuera un campo, cuando en realidad es un método calculado. No se almacena en la
base de datos; se calcula cada vez que se accede.
\end{conceptbox}

\subsection{Order y OrderItem}

\begin{lstlisting}[caption={orders/models.py -- Orden con estados y snapshots}]
class Order(models.Model):
    STATUS_CHOICES = [
        ('pending',   'Pendiente'),
        ('confirmed', 'Confirmado'),
        ('shipped',   'Enviado'),
        ('delivered', 'Entregado'),
        ('cancelled', 'Cancelado'),
    ]
    customer        = models.ForeignKey('customers.Customer',
                          on_delete=models.PROTECT, related_name='orders')
    fecha_creacion  = models.DateTimeField(auto_now_add=True)
    status          = models.CharField(max_length=20, choices=STATUS_CHOICES,
                                       default='pending')
    direccion_envio = models.TextField()
    discount_code   = models.CharField(max_length=50, null=True, blank=True)
    subtotal        = models.DecimalField(max_digits=10, decimal_places=2, ...)
    discount_amount = models.DecimalField(max_digits=10, decimal_places=2, ...)
    total           = models.DecimalField(max_digits=10, decimal_places=2, ...)
    shipping_status = models.CharField(max_length=20, default='pending', ...)
    tracking_number = models.CharField(max_length=100, null=True, blank=True)
\end{lstlisting}

\begin{lstlisting}[caption={orders/models.py -- OrderItem con snapshots del precio}]
class OrderItem(models.Model):
    order        = models.ForeignKey(Order, related_name='items',
                                     on_delete=models.CASCADE)
    variant      = models.ForeignKey('products.ProductVariant',
                                     on_delete=models.PROTECT,
                                     null=True, blank=True)
    product_name = models.CharField(max_length=150)  # snapshot del nombre
    quantity     = models.PositiveIntegerField()
    price        = models.DecimalField(max_digits=10, decimal_places=2, ...)

    @property
    def subtotal(self):
        """Precio unitario x cantidad al momento de la compra."""
        return self.price * self.quantity
\end{lstlisting}

\begin{warnbox}{¿Por qué guardar \texttt{product\_name} y \texttt{price} en el item?}
Si el precio de un producto cambia mañana, las órdenes pasadas deben mantener el
precio que el cliente \textbf{realmente pagó}. Estos campos son \textit{snapshots}
(fotografías en el tiempo). El campo \texttt{variant} es la FK al producto real,
pero \texttt{product\_name} y \texttt{price} son copias inmutables del momento
de la compra.
\end{warnbox}

% ═══════════════════════════════════════════════════════════════
\section{Patrón Builder: Construcción de Objetos Complejos}
% ═══════════════════════════════════════════════════════════════

\begin{conceptbox}{¿Qué problema resuelve el Builder?}
Cuando crear un objeto requiere \textbf{muchos pasos}, validaciones intermedias y
cálculos, un constructor \texttt{\_\_init\_\_} normal se vuelve inmanejable. El
\textbf{Builder} separa la \textit{construcción} del \textit{resultado final}.
Se configura paso a paso y solo al llamar \texttt{build()} se crea el objeto
—ya validado y completo.
\end{conceptbox}

\subsection{OrderBuilder}

\begin{lstlisting}[caption={orders/domain/builders.py -- OrderBuilder completo}]
class OrderBuilder:
    DISCOUNT_RATE = Decimal("0.10")  # 10 % plano

    def __init__(self):
        self._customer       = None
        self._items          = []
        self._direccion_envio = None
        self._discount_code  = None

    # Metodo 1: asignar cliente
    def for_customer(self, customer):
        self._customer = customer
        return self   # <-- retorna self: permite encadenar metodos

    # Metodo 2: agregar un item con validacion de stock
    def add_item(self, variant, quantity: int):
        if quantity <= 0:
            raise ValueError("La cantidad debe ser mayor a 0")

        inventory = variant.inventory
        if not inventory.has_stock(quantity):
            raise ValueError(
                f"Stock insuficiente para '{variant.product.name}'"
            )

        self._items.append({
            "variant":      variant,
            "product_name": variant.product.name,       # snapshot
            "quantity":     quantity,
            "price":        Decimal(str(variant.price)),# snapshot
        })
        return self

    # Metodo 3: asignar direccion de envio
    def with_shipping_address(self, address: str):
        self._direccion_envio = address
        return self

    # Metodo 4: aplicar descuento
    def with_discount(self, discount_code: str = None):
        self._discount_code = discount_code
        return self

    # Metodo final: construir y persistir la orden
    def build(self) -> Order:
        # Validaciones finales obligatorias
        if not self._customer:
            raise ValueError("Se requiere un Customer")
        if not self._items:
            raise ValueError("La orden debe tener al menos un item")
        if not self._direccion_envio:
            raise ValueError("Se requiere direccion de envio")

        # Calculos monetarios
        subtotal = sum(i["quantity"] * i["price"] for i in self._items)
        discount_amount = (subtotal * self.DISCOUNT_RATE
                           if self._discount_code else Decimal("0.00"))
        total = subtotal - discount_amount

        # Persistir la orden
        order = Order.objects.create(
            customer=self._customer,
            direccion_envio=self._direccion_envio,
            discount_code=self._discount_code,
            subtotal=subtotal,
            discount_amount=discount_amount,
            total=total,
        )
        # Persistir cada item
        for item in self._items:
            OrderItem.objects.create(
                order=order,
                variant=item["variant"],
                product_name=item["product_name"],
                quantity=item["quantity"],
                price=item["price"],
            )
        return order
\end{lstlisting}

\begin{tipbox}{Fluent Interface (encadenamiento de métodos)}
Cada método del Builder \texttt{return self} al final. Esto permite escribir:
\begin{verbatim}
orden = (OrderBuilder()
    .for_customer(cliente)
    .with_shipping_address("Calle 123")
    .with_discount("SAVE10")
    .build())
\end{verbatim}
El código es legible como una oración: ``para el cliente, con esta
dirección, con este descuento, construir''.
\end{tipbox}

\subsection{ProductBuilder}

\begin{lstlisting}[caption={products/domain/builders.py -- ProductBuilder}]
class ProductBuilder:
    def __init__(self):
        self._name        = None
        self._description = ''
        self._category    = 'other'
        self._variants    = []

    def with_name(self, name: str):
        if not name or not name.strip():
            raise ValueError("El nombre no puede estar vacio")
        self._name = name.strip()
        return self

    def add_variant(self, sku, price, initial_stock=0,
                    size=None, color=None):
        if not sku or not sku.strip():
            raise ValueError("El SKU no puede estar vacio")
        if price < 0:
            raise ValueError("El precio no puede ser negativo")
        # Validar unicidad del SKU dentro del mismo builder
        if sku in [v['sku'] for v in self._variants]:
            raise ValueError(f"El SKU '{sku}' ya existe")
        self._variants.append({
            "sku": sku, "price": price,
            "initial_stock": initial_stock,
            "size": size, "color": color
        })
        return self

    def build(self):
        if not self._name:
            raise ValueError("El nombre es obligatorio")

        product = Product.objects.create(
            name=self._name,
            description=self._description,
            category=self._category
        )
        for v in self._variants:
            variant = ProductVariant.objects.create(
                product=product, sku=v["sku"],
                price=v["price"], size=v.get("size"), color=v.get("color")
            )
            Inventory.objects.create(
                product_variant=variant,
                available_quantity=v["initial_stock"]
            )
        return product
\end{lstlisting}

% ═══════════════════════════════════════════════════════════════
\section{Patrón Factory: Creación de Dependencias Externas}
% ═══════════════════════════════════════════════════════════════

\begin{conceptbox}{¿Qué problema resuelve el Factory?}
El sistema necesita enviar notificaciones y procesar pagos. En \textbf{desarrollo}
no queremos enviar emails reales ni cobrar dinero real. En \textbf{producción}
sí queremos hacerlo. El problema: \texttt{OrderService} no puede tener un
\texttt{if ENV == 'development'} dentro; eso mezcla responsabilidades.\\[4pt]
La solución: \textbf{una clase Factory} que centraliza esa decisión. El Service
solo pregunta ``dame un notificador'' y la Factory le da el correcto según el entorno.
\textit{El Service nunca sabe si está usando un mock o uno real.}
\end{conceptbox}

\subsection{El árbol de clases completo}

El patrón Factory en este proyecto tiene \textbf{tres niveles}:

\begin{center}
\begin{tikzpicture}[
  cls/.style={draw=accent, fill=boxbg, rounded corners=3pt,
              minimum width=4.2cm, minimum height=0.7cm,
              align=center, font=\small},
  base/.style={draw=warn, fill=warn!10, rounded corners=3pt,
               minimum width=5cm, minimum height=0.8cm,
               align=center, font=\small\bfseries},
  fac/.style={draw=str!70!black, fill=str!10, rounded corners=3pt,
              minimum width=5cm, minimum height=0.8cm,
              align=center, font=\small\bfseries},
  arr/.style={-Stealth, thick},
  inh/.style={-open triangle 60, thick, color=warn}
]
  %% --- Notifiers ---
  \node[base] (bn) at (0,0) {BaseNotifier \\ \scriptsize\texttt{notify(email, msg)}};
  \node[cls]  (mn) at (-2.4,-1.6) {MockNotifier \\ \scriptsize\texttt{[development]}};
  \node[cls]  (en) at ( 2.4,-1.6) {EmailNotifier \\ \scriptsize\texttt{[production]}};
  \draw[inh] (mn) -- (bn);
  \draw[inh] (en) -- (bn);

  \node[fac]  (nf) at (0,-3.2) {NotificationFactory \\ \scriptsize\texttt{create()} $\rightarrow$ BaseNotifier};
  \draw[arr, color=str!70!black] (nf) -- node[left,font=\tiny]{dev} (mn);
  \draw[arr, color=str!70!black] (nf) -- node[right,font=\tiny]{prod} (en);

  %% --- Payment ---
  \node[base] (bp) at (8.5,0) {BasePaymentProcessor \\ \scriptsize\texttt{process\_payment(amount)}};
  \node[cls]  (mp) at (6.1,-1.6) {MockPayment \\ \scriptsize\texttt{[development]}};
  \node[cls]  (rp) at (10.9,-1.6) {RealPayment \\ \scriptsize\texttt{[production]}};
  \draw[inh] (mp) -- (bp);
  \draw[inh] (rp) -- (bp);

  \node[fac]  (pf) at (8.5,-3.2) {PaymentProcessorFactory \\ \scriptsize\texttt{create()} $\rightarrow$ BasePaymentProcessor};
  \draw[arr, color=str!70!black] (pf) -- node[left,font=\tiny]{dev} (mp);
  \draw[arr, color=str!70!black] (pf) -- node[right,font=\tiny]{prod} (rp);
\end{tikzpicture}
\end{center}

\begin{itemize}[noitemsep]
  \item \textbf{Nivel 1 — Base:} Define el \textit{contrato} (qué métodos deben existir).
  \item \textbf{Nivel 2 — Implementaciones:} Dos versiones: Mock (desarrollo) y Real (producción).
  \item \textbf{Nivel 3 — Factory:} Decide cuál instanciar según \texttt{settings.ENV\_TYPE}.
\end{itemize}

\subsection{Nivel 1 — Las Clases Base (el contrato)}

Las clases base definen \textbf{qué operaciones} deben hacer los notificadores y
procesadores, pero \textbf{no las implementan}. Cualquier subclase \textit{está
obligada} a implementar esos métodos.

\begin{lstlisting}[caption={orders/infra/notifiers.py -- BaseNotifier}]
class BaseNotifier:
    """
    Interfaz base para todos los notificadores (LSP).
    Cualquier subclase es intercambiable donde se espere un BaseNotifier.
    """
    def notify(self, user_email: str, message: str) -> dict:
        raise NotImplementedError("Las subclases deben implementar notify()")
\end{lstlisting}

\begin{lstlisting}[caption={orders/infra/payment.py -- BasePaymentProcessor}]
class BasePaymentProcessor:
    """
    Interfaz base para procesadores de pago (LSP).
    Cualquier subclase es intercambiable donde se espere un BasePaymentProcessor.
    """
    def process_payment(self, amount: float) -> dict:
        raise NotImplementedError("Las subclases deben implementar process_payment()")
\end{lstlisting}

\begin{conceptbox}{¿Por qué \texttt{raise NotImplementedError}?}
Si alguien crea una subclase y \textit{olvida} implementar el método, Python
lanzará \texttt{NotImplementedError} en tiempo de ejecución. Es como un
contrato: ``si heredas de mí, \textbf{tienes} que implementar este método o
tu código explota''. Esto garantiza el principio \textbf{LSP} (Liskov):
todo notificador puede usarse donde se espera un \texttt{BaseNotifier}.
\end{conceptbox}

\subsection{Nivel 2 — Las Implementaciones (Mock vs Real)}

\begin{lstlisting}[caption={orders/infra/notifiers.py -- MockNotifier y EmailNotifier}]
class MockNotifier(BaseNotifier):
    # Usado en desarrollo: solo imprime en consola, no envia nada real
    def notify(self, user_email: str, message: str) -> dict:
        print(f"[MOCK NOTIFICATION] To: {user_email}")
        print(f"[MOCK MESSAGE] {message}")
        return {
            'success':   True,
            'method':    'mock',
            'recipient': user_email,
            'message':   'Notification simulated (development mode)'
        }


class EmailNotifier(BaseNotifier):
    # Usado en produccion: aqui iria el codigo real de envio de email
    def notify(self, user_email: str, message: str) -> dict:
        print(f"[REAL EMAIL] Sending to: {user_email}")
        print(f"[EMAIL CONTENT] {message}")
        return {
            'success':   True,
            'method':    'email',
            'recipient': user_email,
            'message':   'Email sent successfully'
        }
\end{lstlisting}

\begin{lstlisting}[caption={orders/infra/payment.py -- MockPaymentProcessor y RealPaymentProcessor}]
class MockPaymentProcessor(BasePaymentProcessor):
    # Genera un ID de transaccion falso, no cobra nada
    def process_payment(self, amount: float) -> dict:
        transaction_id = f"MOCK-{random.randint(100000, 999999)}"
        print(f"[MOCK PAYMENT] Processing ${amount}")
        print(f"[MOCK PAYMENT] Transaction ID: {transaction_id}")
        return {
            'success':        True,
            'transaction_id': transaction_id,
            'amount':         amount,
            'processor':      'mock',
            'message':        'Payment simulated successfully (development mode)'
        }


class RealPaymentProcessor(BasePaymentProcessor):
    # En produccion aqui iria la integracion real con Stripe / PayPal / etc.
    def process_payment(self, amount: float) -> dict:
        print(f"[REAL PAYMENT] Processing ${amount} through payment gateway")
        transaction_id = f"REAL-{random.randint(100000, 999999)}"
        return {
            'success':        True,
            'transaction_id': transaction_id,
            'amount':         amount,
            'processor':      'real',
            'message':        'Payment processed successfully'
        }
\end{lstlisting}

\begin{tipbox}{¿Cuál es la diferencia real entre Mock y Real?}
\begin{itemize}[noitemsep]
  \item \texttt{MockPaymentProcessor}: genera un ID aleatorio con prefijo \texttt{MOCK-}. No llama a ningún banco. Sirve para probar el flujo sin perder dinero.
  \item \texttt{RealPaymentProcessor}: en un sistema real, aquí iría la llamada a la API de Stripe/PayPal/MercadoPago. Genera ID con prefijo \texttt{REAL-}.
\end{itemize}
El \texttt{OrderService} recibe el resultado (un dict con \texttt{success} y
\texttt{transaction\_id}) sin importarle quién lo generó.
\end{tipbox}

\subsection{Nivel 3 — Las Factories (la decisión centralizada)}

\begin{lstlisting}[caption={orders/infra/factories.py -- Las dos factories completas}]
from django.conf import settings
from .notifiers import MockNotifier, EmailNotifier
from .payment   import MockPaymentProcessor, RealPaymentProcessor


class NotificationFactory:

    @staticmethod
    def create():
        # Lee settings.ENV_TYPE = 'development' o 'production'
        if settings.ENV_TYPE == 'development':
            return MockNotifier()   # <-- dev: imprime en consola
        return EmailNotifier()      # <-- prod: envia email real


class PaymentProcessorFactory:

    @staticmethod
    def create():
        if settings.ENV_TYPE == 'development':
            return MockPaymentProcessor()   # <-- dev: simula pago
        return RealPaymentProcessor()       # <-- prod: cobra de verdad
\end{lstlisting}

\begin{conceptbox}{¿Qué hace exactamente \texttt{@staticmethod}?}
Un \texttt{@staticmethod} es un método que \textbf{no depende de instancia}
(\texttt{self}) ni de la clase (\texttt{cls}). Se puede llamar directamente como
\texttt{NotificationFactory.create()} sin necesidad de crear un objeto
\texttt{NotificationFactory()}. Es la forma correcta de un método utilitario
que solo organiza lógica de creación.
\end{conceptbox}

\subsection{¿Cómo usa OrderService las Factories?}

\begin{lstlisting}[caption={orders/services.py -- OrderService usa las Factories}]
class OrderService:

    def __init__(self, notifier=None, payment_processor=None):
        # Si alguien inyecta una dependencia desde afuera (tests), se usa esa.
        # Si no se inyecta nada, la Factory decide cual usar.
        self.notifier = notifier or NotificationFactory.create()
        #                                ^^^^^^^^^^^^^^^^^^^^^
        #                                Retorna MockNotifier o EmailNotifier
        #                                segun ENV_TYPE. OrderService NO sabe cual.

        self.payment_processor = payment_processor or PaymentProcessorFactory.create()
        #                                             ^^^^^^^^^^^^^^^^^^^^^^^^^^^^^^^
        #                                             Retorna Mock o Real.

    def create_order(self, ...):
        # ... mucho codigo ...

        # Llamada al procesador de pago (no importa si es Mock o Real)
        self.payment_processor.process_payment(float(order.total))

        # Llamada al notificador (no importa si es Mock o Email)
        self.notifier.notify(
            user_email=customer.email,
            message=f"Tu orden #{order.pk} fue creada. Total: ${order.total}"
        )
\end{lstlisting}

\begin{warnbox}{El Factory vs el Builder: ¿cuál es la diferencia?}
\begin{itemize}[noitemsep]
  \item \textbf{Builder:} construye un \textit{objeto de dominio} (Order, Product) paso a paso, con validaciones y cálculos acumulados.
  \item \textbf{Factory:} decide \textit{qué implementación} de una dependencia usar, en una sola llamada (\texttt{create()}). No construye nada complejo, solo elige e instancia.
\end{itemize}
En resumen: el Builder \textit{ensambla}, el Factory \textit{elige}.
\end{warnbox}

\begin{tipbox}{Open/Closed Principle en acción}
Si mañana se quiere agregar notificaciones por SMS, el proceso es:
\begin{enumerate}[noitemsep]
  \item Crear \texttt{SMSNotifier(BaseNotifier)} con su propio \texttt{notify()}.
  \item Ajustar \texttt{NotificationFactory.create()} para retornar \texttt{SMSNotifier()}.
  \item \textbf{Nada más.} \texttt{OrderService}, las vistas, los serializers — nada cambia.
\end{enumerate}
Eso es \textit{abierto para extensión, cerrado para modificación (OCP)}.
\end{tipbox}

% ═══════════════════════════════════════════════════════════════
\section{Capa de Aplicación: los Services}
% ═══════════════════════════════════════════════════════════════

Los Services son el \textbf{cerebro del sistema}. Reciben datos del mundo
exterior, orquestan Builders y Factories, y devuelven resultados.
No saben nada de HTTP.

\subsection{OrderService}

\begin{lstlisting}[caption={orders/services.py -- Constructor con inyección de dependencias}]
class StockError(ValueError):
    """Excepcion especifica: la vista la captura y devuelve HTTP 409."""
    pass

class OrderService:
    def __init__(self, notifier=None, payment_processor=None):
        # Si alguien inyecta un notifier (ej: en tests), se usa ese.
        # Si no, la Factory decide cual usar.
        self.notifier          = notifier or NotificationFactory.create()
        self.payment_processor = payment_processor or PaymentProcessorFactory.create()
        self.customer_service  = CustomerService()
        self.product_service   = ProductService()
\end{lstlisting}

\begin{conceptbox}{Dependency Injection (Principio DIP)}
El constructor \texttt{\_\_init\_\_(self, notifier=None)} acepta dependencias
\textbf{desde afuera}. Si alguien pasa un notificador, se usa ese. Si no, la Factory
lo crea. Esto permite en los \textit{tests} hacer:
\begin{verbatim}
service = OrderService(notifier=MockNotifier())
\end{verbatim}
Sin necesidad de modificar la clase para testearla.
\end{conceptbox}

\begin{lstlisting}[caption={orders/services.py -- create\_order (flujo completo)}]
def create_order(self, customer_data, items, shipping_address,
                 discount_code=None) -> dict:
    try:
        # Paso 1: validar que vienen los datos minimos
        if not customer_data or not items or not shipping_address:
            raise ValueError("Faltan datos obligatorios")

        # Paso 2: obtener o crear el cliente
        customer = self.customer_service.get_or_create_customer(
            email=customer_data['email'],
            data=customer_data,
        )

        # Paso 3: iniciar el Builder
        builder = (
            OrderBuilder()
            .for_customer(customer)
            .with_shipping_address(shipping_address)
        )

        # Paso 4: resolver cada variante y añadirla al builder
        for item in items:
            variant = ProductVariant.objects.select_related(
                'product', 'inventory'
            ).get(pk=item['variant_id'])

            try:
                builder.add_item(variant=variant, quantity=item['quantity'])
            except ValueError as stock_err:
                raise StockError(str(stock_err))  # convierte a StockError

        if discount_code:
            builder.with_discount(discount_code)

        # Paso 5: construir y guardar la orden (llama a build())
        order = builder.build()

        # Paso 6: descontar el stock del inventario
        for item in items:
            self.product_service.reserve_stock(
                variant_id=item['variant_id'],
                quantity=item['quantity'],
            )

        # Paso 7: procesar pago (Factory decide Mock vs Real)
        self.payment_processor.process_payment(float(order.total))

        # Paso 8: notificar al cliente (no critico: no detiene el flujo)
        try:
            self.notifier.notify(
                user_email=customer.email,
                message=f"Tu orden #{order.pk} fue creada. Total: ${order.total}"
            )
        except Exception as e:
            print(f"[Warning] Notificacion fallida: {e}")

        return {
            'success':  True,
            'order_id': order.pk,
            'total':    float(order.total),
            'message':  'Orden creada exitosamente',
        }

    except StockError as e:
        return {'success': False, 'error_type': 'stock', 'message': str(e)}
    except ValueError as e:
        return {'success': False, 'message': str(e)}
    except Exception as e:
        return {'success': False, 'message': f'Error: {str(e)}'}
\end{lstlisting}

\begin{conceptbox}{Por qué \texttt{StockError} es una subclase de \texttt{ValueError}}
\texttt{StockError(ValueError)} hereda de \texttt{ValueError}. Esto permite al
\texttt{except StockError} ser más específico que el \texttt{except ValueError}
general. El \textit{Service} sabe distinguir ``falló por stock'' de ``falló por
datos inválidos'', y la \textit{Vista} puede devolver HTTP 409 solo cuando es
un error de stock.
\end{conceptbox}

\subsection{ProductService}

\begin{lstlisting}[caption={products/services.py -- Métodos principales}]
class ProductService:

    def create_product(self, data: dict) -> Product:
        """Delega la creacion al ProductBuilder."""
        builder = ProductBuilder()
        builder.with_name(data['name'])
        builder.with_description(data.get('description', ''))
        builder.with_category(data.get('category', 'other'))
        for v in data.get('variants', []):
            builder.add_variant(sku=v['sku'], price=v['price'],
                                initial_stock=v.get('initial_stock', 0),
                                size=v.get('size'), color=v.get('color'))
        return builder.build()

    def reserve_stock(self, variant_id: int, quantity: int) -> dict:
        """Descuenta stock. Retorna dict con exito/fracaso."""
        variant = ProductVariant.objects.select_related('inventory').get(id=variant_id)
        inventory = variant.inventory
        if not inventory.has_stock(quantity):
            return {'success': False, 'message': 'Stock insuficiente'}
        inventory.available_quantity -= quantity
        inventory.save()
        return {'success': True,
                'message': f'Restante: {inventory.available_quantity}'}

    def update_stock(self, variant_id: int, quantity: int) -> dict:
        """Actualiza el stock absoluto de una variante."""
        inventory = Inventory.objects.get(product_variant__id=variant_id)
        inventory.available_quantity = quantity
        inventory.save()
        return {'success': True, 'new_stock': inventory.available_quantity,
                'is_below_minimum': inventory.is_below_minimum()}
\end{lstlisting}

\subsection{CustomerService}

\begin{lstlisting}[caption={customers/services.py -- get\_or\_create\_customer}]
class CustomerService:

    def get_or_create_customer(self, email: str, data: dict) -> Customer:
        """
        Si ya existe un cliente con ese email: lo devuelve (y actualiza si cambio).
        Si no existe: lo crea.
        """
        customer, created = Customer.objects.get_or_create(
            email=email,
            defaults={
                'nombre':    data.get('nombre', ''),
                'apellido':  data.get('apellido', ''),
                'telefono':  data.get('telefono', ''),
                'direccion': data.get('direccion', ''),
            }
        )
        if not created:  # ya existia: actualizar campos si cambiaron
            for field in ('nombre', 'apellido', 'telefono', 'direccion'):
                if data.get(field):
                    setattr(customer, field, data[field])
            customer.save()
        return customer

    def create_customer(self, data: dict) -> dict:
        """Crea un cliente nuevo. La unicidad de email es logica de negocio:
        vive aqui, NO en el serializer."""
        if Customer.objects.filter(email=data.get('email')).exists():
            return {'success': False, 'customer': None,
                    'message': f"Email {data['email']} ya registrado"}
        customer = Customer.objects.create(**data)
        return {'success': True, 'customer': customer,
                'message': 'Cliente creado exitosamente'}
\end{lstlisting}

\begin{warnbox}{Lógica de negocio en Serializers: ¡PROHIBIDO!}
La unicidad del email es una \textbf{regla de negocio}: ``no pueden existir dos
clientes con el mismo email''. Por eso vive en \texttt{CustomerService}, no en
\texttt{CreateCustomerSerializer}. Los serializers solo validan \textbf{formato}
(¿es un email válido?, ¿el texto no está vacío?), nunca consultan la base de datos.
\end{warnbox}

% ═══════════════════════════════════════════════════════════════
\section{Capa de Presentación: Views y Serializers}
% ═══════════════════════════════════════════════════════════════

\subsection{Serializers: validación de entrada/salida}

\begin{lstlisting}[caption={orders/serializers.py -- CheckoutSerializer (entrada)}]
class CheckoutItemSerializer(serializers.Serializer):
    variant_id   = serializers.IntegerField()       # ID de la variante
    quantity     = serializers.IntegerField(min_value=1)
    # product_name y price son opcionales: el backend los genera como snapshot
    product_name = serializers.CharField(required=False, allow_blank=True)
    price        = serializers.DecimalField(max_digits=10, decimal_places=2,
                                            required=False)

class CheckoutSerializer(serializers.Serializer):
    customer         = CheckoutCustomerSerializer()
    items            = CheckoutItemSerializer(many=True)
    shipping_address = serializers.CharField()
    discount_code    = serializers.CharField(required=False, allow_null=True)

    def validate_items(self, items):
        if not items:
            raise serializers.ValidationError(
                "La orden debe tener al menos un item."
            )
        return items
\end{lstlisting}

\begin{conceptbox}{Interface Segregation Principle en los Serializers}
En lugar de un solo serializer \textit{``HaceloTodo''}, hay serializers específicos:
\begin{itemize}[noitemsep]
  \item \texttt{OrderSerializer} — solo para \textbf{leer} órdenes (GET)
  \item \texttt{CheckoutSerializer} — solo para \textbf{crear} órdenes (POST checkout)
  \item \texttt{CreateProductSerializer} — solo para \textbf{crear} productos
  \item \texttt{UpdateProductSerializer} — solo para \textbf{actualizar} (PATCH)
  \item \texttt{UpdateStockSerializer} — solo para actualizar el stock
\end{itemize}
Cada serializer tiene un propósito único. Eso es \textbf{ISP}.
\end{conceptbox}

\subsection{Views: solo HTTP, sin lógica}

\begin{lstlisting}[caption={orders/views.py -- CreateOrderView (checkout)}]
class CreateOrderView(APIView):
    def post(self, request):
        # 1. Validar el formato de los datos (serializer)
        serializer = CheckoutSerializer(data=request.data)
        if not serializer.is_valid():
            return Response({'errors': serializer.errors},
                            status=status.HTTP_400_BAD_REQUEST)

        # 2. Delegar toda la logica al service
        data   = serializer.validated_data
        result = OrderService().create_order(
            customer_data    = data['customer'],
            items            = data['items'],
            shipping_address = data['shipping_address'],
            discount_code    = data.get('discount_code'),
        )

        # 3. Traducir el resultado del service a un codigo HTTP
        if result['success']:
            return Response(result, status=status.HTTP_201_CREATED)
        if result.get('error_type') == 'stock':
            return Response(result, status=status.HTTP_409_CONFLICT)
        return Response(result, status=status.HTTP_400_BAD_REQUEST)
\end{lstlisting}

\begin{lstlisting}[caption={orders/views.py -- AdminDashboardView}]
class AdminDashboardView(APIView):
    def get(self, request):
        # Una sola linea: delega TODO al service
        stats = OrderService().get_dashboard_stats()
        return Response(stats)
\end{lstlisting}

\begin{tipbox}{La regla de las 15 líneas}
El documento del proyecto dice que una vista con más de 15 líneas probablemente
tiene lógica que debería estar en el Service. Mira \texttt{AdminDashboardView}:
3 líneas. \texttt{CreateOrderView}: 14 líneas incluyendo comentarios.
La vista es solo un \textbf{traductor} entre HTTP y el Service.
\end{tipbox}

% ═══════════════════════════════════════════════════════════════
\section{Resumen de los Principios SOLID en el Proyecto}
% ═══════════════════════════════════════════════════════════════

\begin{center}
\begin{tabular}{@{}lll@{}}
\toprule
\textbf{Principio} & \textbf{Dónde se aplica} & \textbf{Evidencia concreta} \\
\midrule
\textbf{S} — SRP
  & Todas las capas
  & Views $\neq$ Services $\neq$ Builders $\neq$ Factories \\
\textbf{O} — OCP
  & \texttt{NotificationFactory}
  & Agregar \texttt{SMSNotifier} sin tocar \texttt{OrderService} \\
\textbf{L} — LSP
  & \texttt{BaseNotifier}, \texttt{BasePaymentProcessor}
  & Mock y Real son intercambiables \\
\textbf{I} — ISP
  & Serializers
  & \texttt{Create}/\texttt{Update}/\texttt{Read} separados \\
\textbf{D} — DIP
  & \texttt{OrderService.\_\_init\_\_}
  & \texttt{notifier=None} acepta inyección desde fuera \\
\bottomrule
\end{tabular}
\end{center}

% ═══════════════════════════════════════════════════════════════
\section{Flujo Completo de una Compra (Checkout)}
% ═══════════════════════════════════════════════════════════════

Este diagrama muestra el camino que recorre una petición desde el navegador
hasta la base de datos:

\begin{center}
\begin{tikzpicture}[
  box/.style={draw=accent, rounded corners=3pt, fill=boxbg,
              text width=3.2cm, align=center, font=\small,
              minimum height=0.8cm},
  arrow/.style={-Stealth, thick, color=accent},
  label/.style={font=\tiny, color=cmt}
]
  \node[box] (browser)  {Navegador\\(React)};
  \node[box, right=1.2cm of browser] (view)    {CreateOrderView\\(View)};
  \node[box, right=1.2cm of view]    (ser)     {CheckoutSerializer\\(Serializer)};
  \node[box, below=1cm of view]      (svc)     {OrderService\\(Service)};
  \node[box, below=1cm of svc]       (builder) {OrderBuilder\\(Builder)};
  \node[box, right=1.2cm of svc]     (db)      {Base de Datos\\(Django ORM)};
  \node[box, left=1.2cm of svc]      (factory) {Factories\\(Notif + Pay)};

  \draw[arrow] (browser) -- node[above,label]{POST /checkout} (view);
  \draw[arrow] (view) -- node[above,label]{valida} (ser);
  \draw[arrow] (ser)  -- node[right,label]{datos válidos} (view);
  \draw[arrow] (view) -- node[right,label]{delega} (svc);
  \draw[arrow] (svc)  -- node[right,label]{construye} (builder);
  \draw[arrow] (builder) -- node[above,label]{crea registros} (db);
  \draw[arrow] (svc) -- node[above,label]{notifica/paga} (factory);
  \draw[arrow] (factory) -- node[left,label]{Mock en dev} (db);
  \draw[arrow] (svc) -- node[above,label]{201 OK} (view);
  \draw[arrow] (view) -- node[above,label]{Response} (browser);
\end{tikzpicture}
\end{center}

\vspace{1em}
\noindent\textbf{Paso a paso:}

\begin{enumerate}
  \item El \textbf{navegador} envía \texttt{POST /api/orders/checkout/} con datos JSON.
  \item La \textbf{vista} pasa los datos al \texttt{CheckoutSerializer}.
  \item El \textbf{serializer} valida formato (campos obligatorios, tipos). Si hay error → HTTP 400.
  \item La \textbf{vista} delega al \texttt{OrderService.create\_order()}.
  \item El \textbf{service} llama a \texttt{CustomerService} para obtener/crear el cliente.
  \item El \textbf{service} inicia el \texttt{OrderBuilder} y le agrega variantes una a una.
  \item El \textbf{builder} valida el stock en cada \texttt{add\_item()}.
  \item El \textbf{builder} calcula totales y llama a \texttt{Order.objects.create()}.
  \item El \textbf{service} llama a \texttt{ProductService.reserve\_stock()} para cada ítem.
  \item El \textbf{service} llama a \texttt{PaymentProcessorFactory.create().process\_payment()}.
  \item El \textbf{service} llama a \texttt{NotificationFactory.create().notify()}.
  \item La \textbf{vista} devuelve HTTP 201 con el ID y total de la orden.
\end{enumerate}

% ═══════════════════════════════════════════════════════════════
\section{Estructura de Carpetas y su Justificación}
% ═══════════════════════════════════════════════════════════════

\begin{verbatim}
CEFUSAECommerce/
 ├── products/          # App independiente: Product, Variant, Inventory
 │   ├── models.py
 │   ├── services.py    # ProductService
 │   ├── serializers.py
 │   ├── views.py
 │   └── domain/
 │       └── builders.py  # ProductBuilder
 ├── customers/         # App independiente: Customer
 │   ├── models.py
 │   ├── services.py    # CustomerService
 │   ├── serializers.py
 │   └── views.py
 └── orders/            # App independiente: Order, OrderItem
     ├── models.py
     ├── services.py    # OrderService (orquesta todo)
     ├── serializers.py
     ├── views.py
     ├── domain/
     │   └── builders.py   # OrderBuilder (patron Builder)
     └── infra/
         ├── factories.py  # NotificationFactory, PaymentProcessorFactory
         ├── notifiers.py  # BaseNotifier, MockNotifier, EmailNotifier
         └── payment.py    # BasePaymentProcessor, Mock, Real
\end{verbatim}

\begin{tipbox}{¿Por qué \texttt{domain/} e \texttt{infra/} son subcarpetas de \texttt{orders/}?}
\begin{itemize}[noitemsep]
  \item \texttt{domain/} — contiene lógica \textit{pura de negocio} (Builders). No depende de HTTP ni de servicios externos.
  \item \texttt{infra/} — contiene \textit{detalles técnicos} (Factories, notificadores, pasarelas de pago). Pueden cambiar sin afectar la lógica de negocio.
\end{itemize}
Esta separación es la base de la \textbf{Arquitectura Limpia} (Clean Architecture).
\end{tipbox}

% ═══════════════════════════════════════════════════════════════
\section{Endpoints Disponibles}
% ═══════════════════════════════════════════════════════════════

\begin{center}
\begin{tabular}{@{}lll@{}}
\toprule
\textbf{Método + URL} & \textbf{Descripción} & \textbf{HTTP} \\
\midrule
\texttt{POST /api/orders/checkout/} & Crear una orden (checkout) & 201 / 400 / 409 \\
\texttt{GET /api/orders/<id>/}      & Detalle de una orden & 200 / 404 \\
\texttt{PATCH /api/orders/<id>/status/} & Cambiar estado & 200 / 400 \\
\midrule
\texttt{GET /api/products/}         & Listar productos activos & 200 \\
\texttt{POST /api/products/}        & Crear producto & 201 / 400 \\
\texttt{GET /api/products/<id>/}    & Detalle de producto & 200 / 404 \\
\texttt{PATCH /api/products/<id>/}  & Actualizar producto & 200 / 404 \\
\texttt{DELETE /api/products/<id>/} & Desactivar producto & 204 / 404 \\
\texttt{POST /api/products/check-stock/} & Verificar stock & 200 \\
\midrule
\texttt{GET /api/customers/}        & Listar clientes & 200 \\
\texttt{POST /api/customers/}       & Crear cliente & 201 / 409 \\
\texttt{GET /api/customers/<id>/}   & Detalle de cliente & 200 / 404 \\
\texttt{GET /api/customers/<id>/orders/} & Órdenes del cliente & 200 / 404 \\
\midrule
\texttt{GET /api/admin/dashboard/}  & Estadísticas admin & 200 \\
\texttt{GET /api/admin/orders/}     & Listar órdenes (admin) & 200 \\
\texttt{DELETE /api/admin/orders/<id>/} & Eliminar orden & 204 / 404 \\
\bottomrule
\end{tabular}
\end{center}

\end{document}
